%% lp-core-petta.tex: Certified PeTTa — LP Semantics, Prolog Compilation, and OSLF Integration
%%
%% Render with:
%%   pdflatex lp-core-petta.tex && bibtex lp-core-petta && pdflatex lp-core-petta.tex && pdflatex lp-core-petta.tex

\documentclass[]{article}
\usepackage{url}
\usepackage[hyperindex,breaklinks]{hyperref}
\usepackage{listings}
\lstset{basicstyle=\ttfamily\footnotesize,breaklines=true,frame=single}
\usepackage{float}
\restylefloat{table}
\usepackage{longtable}
\usepackage{graphicx}
\usepackage[font=small,labelfont=bf]{caption}
\usepackage[skip=0pt]{subcaption}
\usepackage{amsfonts}
\usepackage{amssymb}
\usepackage{amsmath}
\usepackage{amsthm}
\usepackage{tabularx}
\usepackage{array}

\newtheorem{theorem}{Theorem}
\newtheorem{lemma}[theorem]{Lemma}
\newtheorem{definition}[theorem]{Definition}
\newtheorem{corollary}[theorem]{Corollary}

\newcommand{\code}[1]{\texttt{#1}}
\newcommand{\codepath}[1]{\path{#1}}
\newcommand{\MeTTa}{\textup{MeTTa}}
\newcommand{\PeTTa}{\textup{PeTTa}}
\newcommand{\MeTTaIL}{\textup{MeTTaIL}}
\newcommand{\OSLF}{\textup{OSLF}}
\newcolumntype{Y}{>{\raggedright\arraybackslash}X}
\setlength{\emergencystretch}{3em}
\sloppy

\lstdefinelanguage{lean}{
  keywords={def,theorem,lemma,structure,inductive,where,by,exact,simp,intro,
             cases,match,with,fun,let,have,show,rfl,and,or,not,if,then,else,
             open,namespace,end,abbrev,instance,class,return,do,for,in,
             import,private,protected,noncomputable,mutual,termination_by},
  sensitive=true,
  comment=[l]{--},
  morecomment=[s]{/-}{-/},
  morestring=[b]",
  literate={->}{{$\to$}}1 {<-}{{$\leftarrow$}}1
           {forall}{{$\forall$}}1
           {exists}{{$\exists$}}1
           {<=}{{$\le$}}1 {>=}{{$\ge$}}1
}

\title{Certified \PeTTa{}:\\Logic Programming Semantics, Prolog Compilation,\\and \OSLF{} Integration in Lean~4}
\author{Zar \and Oru\v{z}i\footnote{``Oru\v{z}i'' denotes AI-assisted collaboration using Claude Code (Anthropic), ChatGPT, and Codex (OpenAI).  Lean~4 formalization at \url{https://github.com/mettapedia}.}}

\begin{document}
\maketitle

\begin{abstract}
We present a Lean~4 formalization of the \PeTTa{} language---the Prolog-based
implementation of \MeTTa{}---comprising 13{,}956 lines across 44 files with zero
sorries and zero custom axioms.  The formalization establishes a certified
semantic stack from LP foundations through Prolog-style goal evaluation to an
\OSLF{} type system with Galois connection and Rust code export.

The main contributions are: (1)~a full LP kernel with $T_P$, least Herbrand
model via Tarski's theorem, SLD resolution, and assumption-free unification
completeness; (2)~correctness of \MeTTaIL{} pattern matching on the
\code{isMatchCorrect} fragment; (3)~the \MeTTaIL{}--LP bridge proving that
\code{DeclReduces} steps imply least-model membership; (4)~\code{PeTTaEval},
a pure evaluation relation with LP soundness for rule application;
(5)~\code{PeTTaCmd}, a stateful effects layer for \code{add-atom},
\code{remove-atom}, \code{get-atoms}, \code{progn}, and \code{prog1};
(6)~a certified propositional connection-tableau chainer with DFS
refinement and witness-level completeness; (7)~a Prolog-style goal
language (\code{PrologGoal}) with conjunctive/disjunctive goals, \code{findall},
and \code{assert}, together with a compilation pass
(\code{compileExpr}) that translates \PeTTa{} expressions into Prolog goals with
correctness theorems; (8)~a 4-argument \MeTTa{} evaluation judgment
(\code{MeTTaEval}) with grounded oracles, standard library derived forms, and
typed evaluation; and (9)~an \OSLF{} pipeline instance composing
\code{pettaSpaceToLangDef} with the generic \OSLF{} type synthesis to produce
a type system with automatic $\Diamond \dashv \Box$ Galois connection and Rust
\code{language!\{\}} macro export, plus a GSLT forward fiber for categorical
integration.

Together these establish a certified end-to-end pipeline from \PeTTa{} source
semantics through LP model theory to a deployable \OSLF{} type system.
\end{abstract}

\tableofcontents

%% ====================================================================
\section{Introduction}
%% ====================================================================

\MeTTa{} (Meta Type Talk) is the native programming language of the OpenCog Hyperon
AGI framework~\cite{goertzel2009pln}.  It combines higher-order term rewriting
with probabilistic inference and a first-class atomspace (knowledge graph), making
it a technically demanding target for formal verification.  \PeTTa{} is the
operational, Prolog-based implementation of \MeTTa{}, implemented as a transpiler
from \MeTTa{} surface syntax to SWI-Prolog.

Our goal in this paper is to give the pure, rule-driven fragment of \PeTTa{} a
certified semantic foundation.  Logic programs with definite clauses provide exactly
the right vehicle: the least Herbrand model semantics~\cite{vanEmdenKowalski1976,Lloyd1987}
is compositional, fixed-point-theoretic, and mechanically tractable in Lean~4.

\paragraph{Contribution overview.}
We develop and verify the following stack in Lean~4 (all with zero sorries and zero
custom axioms):

\begin{enumerate}
\item \textbf{LP kernel} (Sections~\ref{sec:lp-kernel}--\ref{sec:bridge}).
  $T_P$ operator, least Herbrand model via Tarski's theorem, SLD resolution,
  assumption-free unification completeness, and the \MeTTaIL{}--LP bridge
  proving that \code{DeclReduces} steps imply least-model membership.

\item \textbf{\MeTTaIL{} pattern matching} (Section~\ref{sec:matchspec}).
  Correctness of the algorithmic matcher on the \code{isMatchCorrect} fragment.

\item \textbf{\PeTTa{} evaluation and LP soundness} (Section~\ref{sec:petta}).
  The \code{PeTTaEval} relation captures pure evaluation; LP soundness reflects
  every \code{ruleApp} step in the least Herbrand model.

\item \textbf{Stateful effects} (Section~\ref{sec:effects}).
  \code{PeTTaCmd} models atomspace mutations with state monotonicity.

\item \textbf{Certified ATP/chainer} (Section~\ref{sec:chainer}).
  Connection-tableau DFS with soundness and witness-level completeness.

\item \textbf{Prolog goal language and compilation} (Section~\ref{sec:prolog}).
  A \code{PrologGoal} type with conjunction, disjunction, \code{findall}, and
  \code{assert}; a compilation pass \code{compileExpr} from \PeTTa{} expressions
  to Prolog goals with correctness theorems.

\item \textbf{4-argument \MeTTa{} evaluation} (Section~\ref{sec:mettaeval}).
  \code{MeTTaEval s ctx expr typ} with grounded oracles, standard library derived
  forms (\code{if-equal}, \code{car-atom}, \code{let}), and type-driven evaluation.

\item \textbf{\OSLF{} pipeline and GSLT integration} (Section~\ref{sec:oslf}).
  \code{pettaOSLF} composes the \PeTTa{} space into the generic \OSLF{} type
  synthesis, obtaining an automatic $\Diamond \dashv \Box$ Galois connection.
  Rust \code{language!\{\}} macro export and a unit-indexed GSLT forward fiber
  complete the integration.
\end{enumerate}

\paragraph{Architecture diagram.}
The certified pipeline:

\begin{center}
\fbox{\parbox{0.90\linewidth}{
\texttt{pettaOSLF} $\longrightarrow$ \OSLF{} type system ($\Diamond \dashv \Box$)
  $\longrightarrow$ \texttt{pettaRenderRust}\\
\texttt{pettaForwardFiber} (GSLT categorical integration)\\[2pt]
$\uparrow$ \texttt{pettaSpaceToLangDef}\\[2pt]
\texttt{MeTTaEval} (4-arg: space, context, expr, type)\\
$\downarrow$ \texttt{compileExpr} (PeTTa $\to$ PrologGoal)\\
\texttt{PeTTaCmd} (stateful: add/remove/get/progn/prog1)\\
$\downarrow$ pureEval lift\\
\texttt{PeTTaEval} (pure: var/ruleApp/spaceQuery/superpose/collapse)\\
$\downarrow$ \texttt{matchPattern\_correct} + \texttt{pettaRuleSafe}\\
$\downarrow$ \texttt{lp\_complete\_topRule} (\texttt{MeTTaILBridge})\\
$\downarrow$ \texttt{leastHerbrandModel} ($T_P$ fixed point)
}}
\end{center}

\paragraph{Paper structure.}
Sections~\ref{sec:background}--\ref{sec:bridge} cover background, the LP kernel,
\MeTTaIL{} pattern matching, and the \MeTTaIL{}--LP bridge.
Section~\ref{sec:petta} presents \code{PeTTaEval} and LP soundness.
Section~\ref{sec:chainer} adds the certified ATP/chainer layer.
Section~\ref{sec:effects} covers the stateful effects layer.
Section~\ref{sec:typesystem} formalises the \MeTTa{} type system fragment.
Section~\ref{sec:prolog} presents the Prolog goal language and \code{compileExpr}.
Section~\ref{sec:mettaeval} covers the 4-argument \code{MeTTaEval} judgment,
grounded oracles, and standard library forms.
Section~\ref{sec:oslf} integrates the \OSLF{} pipeline and GSLT fiber.
Section~\ref{sec:related} discusses related work and Section~\ref{sec:conclusion}
concludes.

%% ====================================================================
\section{Background}
\label{sec:background}
%% ====================================================================

\paragraph{Logic programming.}
A \emph{definite clause program} (or LP) is a finite set of \emph{Horn clauses}
$H \leftarrow B_1, \ldots, B_n$ where $H$ and each $B_i$ are atoms over a fixed
signature.  The \emph{Herbrand base} is the set of all ground atoms constructible
from the function and constant symbols in the signature.  The immediate consequence
operator $T_P$ maps an interpretation $I$ (a subset of the Herbrand base) to the
set of all ground atoms that are immediate consequences of applying a grounded
clause whose body is satisfied by $I$.  Since $T_P$ is monotone on the complete
lattice $(\mathcal{P}(\text{HB}), \subseteq)$, the Knaster--Tarski theorem
guarantees a least fixed point $T_P \uparrow \omega$, which equals the least
Herbrand model~\cite{vanEmdenKowalski1976,Lloyd1987}.

\paragraph{\MeTTa{} and \PeTTa{}.}
\MeTTa{} expressions are terms over a \emph{type-free} base (atoms, applications,
variables).  The semantics of a \MeTTa{} program is given by its rewrite rules
\code{(= lhs rhs)}: an expression reduces to all terms reachable by matching
its subexpressions against LHS patterns and substituting the corresponding RHS.
The non-determinism is explicit: \code{superpose} collects alternatives and
\code{collapse} gathers all answers via \code{findall}-style comprehension.

\PeTTa{}~\cite{petta2024} is the Prolog-based implementation of \MeTTa{}.
The transpiler (\code{transpiler.pl}) converts \MeTTa{} rules into Prolog clauses;
\code{match \&self} queries are handled by \code{match\_term/3} in \code{spaces.pl};
\code{superpose}/\code{collapse} map to Prolog disjunction and \code{findall/3}.

\paragraph{Lean~4.}
Lean~4~\cite{Moura2021} is a dependently-typed theorem prover and programming
language.  Our formalization uses Mathlib~\cite{mathlib2020} for Tarski's fixed-point
theorem (\code{OrderHom.lfp}) and basic order theory.  The \code{@[simp]} attribute
marks rewrite lemmas for the \code{simp} tactic; \code{decide} proves decidable
propositions by kernel reduction.

%% ====================================================================
\section{The LP Kernel}
\label{sec:lp-kernel}
%% ====================================================================

The LP kernel lives in \code{Mettapedia/Logic/LP/} and spans three files:
\code{Substitution.lean} (206 lines), \code{Semantics.lean} (204 lines), and
\code{SLD.lean} (169 lines), totalling 579 lines.

\subsection{Signature and Terms}

\begin{lstlisting}[language=lean]
structure LPSignature where
  vars            : Type
  constants       : Type
  functionSymbols : Type
  functionArity   : functionSymbols -> Nat
  relationSymbols : Type
  relationArity   : relationSymbols -> Nat

inductive Term (sigma : LPSignature) where
  | var   : sigma.vars -> Term sigma
  | const : sigma.constants -> Term sigma
  | app   : (f : sigma.functionSymbols) ->
            (Fin (sigma.functionArity f) -> Term sigma) -> Term sigma

structure Atom (sigma : LPSignature) where
  symbol : sigma.relationSymbols
  args   : Fin (sigma.relationArity symbol) -> Term sigma
\end{lstlisting}

Ground terms and atoms are defined analogously with \code{GroundTerm} and
\code{GroundAtom}, using the same finitely-typed argument representation.  A
\emph{grounding} \code{Grounding sigma := sigma.vars -> GroundTerm sigma}
maps each LP variable to a ground term; applying it to a \code{Term} or \code{Atom}
yields the corresponding \code{GroundTerm} or \code{GroundAtom}.

\subsection{Clauses and Knowledge Bases}

\begin{lstlisting}[language=lean]
structure Clause (sigma : LPSignature) where
  head : Atom sigma
  body : List (Atom sigma)

structure GroundClause (sigma : LPSignature) where
  head : GroundAtom sigma
  body : List (GroundAtom sigma)

abbrev Program (sigma)       := List (Clause sigma)
abbrev Interpretation (sigma) := Set (GroundAtom sigma)

structure KnowledgeBase (sigma : LPSignature) where
  prog : Program sigma
  db   : Set (GroundAtom sigma)   -- extensional (EDB) ground atoms
\end{lstlisting}

The \code{db} field holds extensional (EDB) ground facts; \code{prog} holds
intensional (IDB) definite clauses.

\subsection{Least Herbrand Semantics}

The immediate consequence operator:

\begin{lstlisting}[language=lean]
noncomputable def T_P_LP {sigma : LPSignature}
    (kb : KnowledgeBase sigma) (I : Interpretation sigma) :
    Interpretation sigma :=
  kb.db
  union { a | exists (c : Clause sigma) (g : Grounding sigma),
              c in kb.prog
              /\ g.groundAtom c.head = a
              /\ forall b in c.body, g.groundAtom b in I }
\end{lstlisting}

Monotonicity of $T_P$:

\begin{theorem}[Monotonicity]
$T_P$ is monotone: if $I \subseteq J$ then $T_P(I) \subseteq T_P(J)$.
\end{theorem}

\begin{proof}
Straightforward: for each clause witness $(c, g)$ satisfying the body in $I$,
the same witness satisfies the body in $J \supseteq I$.
\end{proof}

Packaging $T_P$ as an \code{OrderHom} (an order-homomorphism on the complete
lattice of sets), Tarski's least-fixed-point theorem gives the least Herbrand model:

\begin{lstlisting}[language=lean]
noncomputable def leastHerbrandModel {sigma : LPSignature}
    (kb : KnowledgeBase sigma) : Interpretation sigma :=
  OrderHom.lfp (T_P_LP_orderHom kb)
\end{lstlisting}

The following three theorems are the cornerstones of the LP semantics:

\begin{theorem}[\code{leastHerbrandModel\_fixpoint}]
\label{thm:lfp}
$T_P(\mathrm{LHM}(kb)) = \mathrm{LHM}(kb)$.
\end{theorem}

\begin{theorem}[\code{leastHerbrandModel\_clause}]
\label{thm:clause}
If $c \in kb.prog$, $g$ is a grounding, and for all $b \in c.body$,
$g(b) \in \mathrm{LHM}(kb)$, then $g(c.head) \in \mathrm{LHM}(kb)$.
\end{theorem}

\begin{theorem}[\code{leastHerbrandModel\_least}]
\label{thm:least}
If $T_P(I) \subseteq I$, then $\mathrm{LHM}(kb) \subseteq I$.
\end{theorem}

The proofs are one-liners using \code{OrderHom.isFixedPt\_lfp} and \code{OrderHom.lfp\_le}
from Mathlib.  Theorem~\ref{thm:clause} follows immediately from
Theorem~\ref{thm:lfp} and the definition of $T_P$: if the body holds in
the model (which is a fixed point), then the head holds by one application of $T_P$.

\subsection{SLD Resolution and All-Solutions SLD}

The file \code{SLD.lean} (169 lines) formalises SLD resolution as an inductive type:

\begin{lstlisting}[language=lean]
inductive SLDTree (sigma : LPSignature) (kb : KnowledgeBase sigma)
    : List (Atom sigma) -> Subst sigma -> Prop where
  | success : SLDTree sigma kb [] theta
  | resolve : ...
\end{lstlisting}

The soundness theorem (\code{sldTree\_sound}) states that if a tree witnesses a
successful derivation of a goal $G$ with substitution $\theta$, then the grounded
goal $\theta(G)$ is in the least Herbrand model.  The file \code{SLDAll.lean}
(205 lines) provides \code{sldSearchAll}, which collects all answer substitutions,
together with its soundness theorem.

\subsection{Operational Completeness: Unification and SLD}

The LP completeness layer extends the kernel with explicit executable-success
theorems.

\paragraph{Assumption-free semantic completeness of unification.}
\code{Logic/LP/UnificationComplete.lean} (513 lines) internalizes the
semantic-to-derivation bridge and proves:

\begin{theorem}[\code{unifyFuel\_exists\_of\_unifies}]
\label{thm:unifycomplete}
If an equation list is semantically unifiable, then \code{unifyFuel} succeeds
for some finite fuel:
\[
  (\exists \delta,\ \code{Unifies}\ \delta\ \code{eqs}) \Rightarrow
  (\exists fuel\ \theta,\ \code{unifyFuel}\ fuel\ \code{eqs} = \code{some}\ \theta).
\]
\end{theorem}

Function-free and ground specializations are also proved:
\begin{itemize}
\item \codepath{unifyFuel_exists_of_unifies_functionFree}
\item \codepath{unifyFuel_exists_of_unifies_ground}
\item \codepath{unifyFuel_exists_of_unifies_functionFree_ground}
\end{itemize}
with canary theorems for occurs-check rejection and non-occurs success in
\codepath{UnificationCompletenessCanaries.lean}.

\paragraph{SLD completeness bundles from least-model lifting.}
\code{Logic/LP/SLDCompute.lean} (485 lines) introduces \code{SLDWitness}
and proves that witness existence implies executable success
(\code{sldSearch\_complete\_bounded}, \code{sldQuery\_complete\_bounded}).
The main endpoint is:

\begin{theorem}[\code{sldQuery\_complete\_of\_semantic\_lift}]
\label{thm:sldcomplete}
Given lifting data from least-model facts and one-step rule closures to
\code{SLDWitness}, every least-Herbrand-model ground atom has a successful
executable query:
\[
  ga \in \mathrm{LHM}(kb) \Rightarrow
  \exists fuel\ \theta,\ \code{sldQuery}\ kb.prog\ ga\ fuel = \code{some}\ \theta.
\]
\end{theorem}

The bundled endpoint
\codepath{sldQuery_complete_sound_bundle_of_semantic_lift}
adds soundness of returned substitutions in the same theorem. Canonical
kit instances in \code{SLDCompletenessKit.lean} (261 lines) provide reusable
hypothesis packages for concrete KB classes (e.g., nullary fact-only and
nullary unary-body fragments).

%% ====================================================================
\section{\MeTTaIL{} Pattern Matching Correctness}
\label{sec:matchspec}
%% ====================================================================

\code{OSLF/MeTTaIL/Match.lean} (201 lines) implements the algorithmic pattern
matcher; \code{OSLF/MeTTaIL/MatchSpec.lean} (821 lines) proves its correctness.

\subsection{\MeTTaIL{} Syntax}

\MeTTaIL{} uses a \emph{locally nameless} representation~\cite{AydemarPOPL2008}:
$\alpha$-equivalent patterns are syntactically identical, so no renaming is needed
during matching or substitution.

\begin{lstlisting}[language=lean]
inductive Pattern where
  | fvar        : String -> Pattern       -- free variable (metavariable)
  | bvar        : Nat -> Pattern          -- de Bruijn bound variable
  | apply       : String -> List Pattern -> Pattern
  | lambda      : Pattern -> Pattern
  | multiLambda : Nat -> Pattern -> Pattern
  | collection  : CollType -> List Pattern -> Option String -> Pattern
  | subst       : Pattern -> Pattern -> Pattern   -- explicit substitution

inductive CollType where | vec | hashBag | hashSet

structure RewriteRule where
  left     : Pattern
  right    : Pattern
  premises : List Pattern   -- [] for unconditional rules
\end{lstlisting}

The \code{fvar} constructor represents \emph{pattern metavariables}: they match any
term and produce a variable binding.  The \code{bvar n} constructor is a de Bruijn
index for bound variables; it matches only \code{bvar m} with $m = n$.
The \code{.subst body repl} form encodes explicit $\beta$-reduction:
\code{applyBindings} evaluates it by calling \code{openBVar 0 repl body}.

\subsection{Algorithmic Matching}

\begin{lstlisting}[language=lean]
def matchPattern (pat term : Pattern) : List Bindings :=
  match pat, term with
  | .fvar x, t                          => [[(x, t)]]
  | .bvar n, .bvar m                    => if n == m then [[]] else []
  | .apply c1 pargs, .apply c2 targs   =>
      if c1 == c2 && pargs.length == targs.length
      then matchArgs pargs targs
      else []
  | .lambda bodyPat, .lambda bodyConcrete =>
      matchPattern bodyPat bodyConcrete
  | .multiLambda npat bodyPat, .multiLambda nconc bodyConcrete =>
      if npat == nconc then matchPattern bodyPat bodyConcrete else []
  | .collection ct1 pelems rest1, .collection ct2 telems _rest2 =>
      if ct1 == ct2 then matchBag pelems rest1 ct1 telems else []
  | .subst pbody prepl, .subst tbody trepl =>
      (matchPattern pbody tbody).flatMap fun b1 =>
        (matchPattern prepl trepl).filterMap (mergeBindings b1)
  | _, _                                => []
\end{lstlisting}

The function is mutually recursive with \code{matchArgs} (pairwise argument
matching) and \code{matchBag} (multiset matching).  Termination is by
\code{sizeOf pat}.  The return type \code{List Bindings} encodes non-determinism:
bag matching may succeed in multiple ways.

\code{applyBindings} substitutes \code{fvar} metavariables by their bindings and
evaluates \code{.subst} nodes:

\begin{lstlisting}[language=lean]
def applyBindings (bindings : Bindings) (rhs : Pattern) : Pattern :=
  match rhs with
  | .fvar x   => match bindings.find? (fun p => p.1 == x) with
                 | some (_, val) => val
                 | none          => .fvar x
  | .bvar n   => .bvar n
  | .apply c args => .apply c (args.map (applyBindings bindings))
  | .lambda body  => .lambda (applyBindings bindings body)
  | .multiLambda n body => .multiLambda n (applyBindings bindings body)
  | .subst body repl =>
      openBVar 0 (applyBindings bindings repl)
                 (applyBindings bindings body)
  | .collection ct elems rest =>
      let elems' := elems.map (applyBindings bindings)
      let restElems := match rest with
        | some rv => match bindings.find? (fun p => p.1 == rv) with
                     | some (_, .collection _ relems _) => relems
                     | _ => []
        | none => []
      .collection ct (elems' ++ restElems) none
\end{lstlisting}

\subsection{The \texttt{isMatchCorrect} Fragment}

Not every pattern $p$ satisfies $\code{applyBindings}(bs, p) = t$ for all
$bs \in \code{matchPattern}(p, t)$.  The two problematic constructors are:
\begin{itemize}
\item \code{.collection}: bag matching picks elements out of order, so the
  reconstituted pattern may have a different element permutation than $t$.
\item \code{.subst}: \code{applyBindings} evaluates the substitution node via
  \code{openBVar}, which is not the identity.
\end{itemize}

The Boolean predicate \code{Pattern.isMatchCorrect} excludes both:

\begin{lstlisting}[language=lean]
def Pattern.isMatchCorrect (p : Pattern) : Bool :=
  match p with
  | .fvar _           => true
  | .bvar _           => true
  | .apply _ args     => args.all isMatchCorrect
  | .lambda body      => isMatchCorrect body
  | .multiLambda _ b  => isMatchCorrect b
  | .subst _ _        => false
  | .collection _ _ _ => false
\end{lstlisting}

\subsection{Correctness Theorem}

The proof architecture in \code{MatchSpec.lean} is:
(1)~define relational specifications \code{MatchRel}, \code{MatchArgsRel},
\code{MatchBagRel} as mutual inductives independent of the algorithm;
(2)~prove algorithmic soundness (if $bs \in \code{matchPattern}(p, t)$ then
$\code{MatchRel}(p, t, bs)$) by strong induction on \code{sizeOf pat};
(3)~prove that \code{MatchRel} implies \code{applyBindings bs pat = t} on the
\code{isMatchCorrect} fragment (by structural induction on \code{MatchRel}).

\begin{theorem}[\code{matchPattern\_correct}]
\label{thm:matchcorrect}
If $bs \in \code{matchPattern}(pat, t)$ and $\code{pat.isMatchCorrect} = \code{true}$,
then $\code{applyBindings}(bs, pat) = t$.
\end{theorem}

\begin{proof}
By the soundness bridge, $bs \in \code{matchPattern}(pat, t)$ implies
$\code{MatchRel}(pat, t, bs)$.  Then \code{matchRel\_correct} (structural induction
on \code{MatchRel}, using the \code{isMatchCorrect} guard to exclude collection and
subst cases) gives $\code{applyBindings}(bs, pat) = t$.
\end{proof}

The \code{isMatchCorrect} guard is sharp:

\begin{theorem}[\code{matchPattern\_correct\_false}]
There exist a pattern $pat$, a term $t$, and bindings $bs$
such that $bs \in \code{matchPattern}(pat, t)$ yet
$\code{applyBindings}(bs, pat) \neq t$.
\end{theorem}

\begin{proof}
Counterexample:
\[
  pat = \mathrm{vec}(x,y),\quad t = \mathrm{vec}(b,a).
\]
Bag matching with $i = 1$ first yields
$bs = [(\text{"y"}, b), (\text{"x"}, a)]$, so
\[
  \code{applyBindings}(bs, pat) = [a,b] \neq [b,a] = t.
\]
\end{proof}

%% ====================================================================
\section{The \MeTTaIL{}--LP Bridge}
\label{sec:bridge}
%% ====================================================================

The bridge lives in \code{Logic/LP/MeTTaILBridge.lean} (783 lines).

\subsection{LP Signature for \MeTTaIL{}}

\begin{lstlisting}[language=lean]
abbrev mettailLPSig : LPSignature where
  vars            := String
  constants       := String
  functionSymbols := Unit      -- single binary pairing function
  functionArity   := fun _ => 2
  relationSymbols := Unit      -- single "reduces" relation
  relationArity   := fun _ => 2
\end{lstlisting}

Using \code{abbrev} (transparent) rather than \code{def} (opaque) is essential:
it makes \code{mettailLPSig.functionArity () = 2} definitionally true, so that
the \code{Matrix.cons} expression \code{![a, b]} typechecks for
\code{Fin (mettailLPSig.functionArity ()) -> T}.

Patterns are encoded as binary pair-spine trees:

\begin{lstlisting}[language=lean]
def pairTerm (a b : Term mettailLPSig) : Term mettailLPSig :=
  Term.app () ![a, b]

-- Encoding:
--   .fvar x     ->  Term.var x              (LP variable)
--   .bvar n     ->  Term.const "bv:n"       (ground token)
--   .apply f as ->  pair("f_f", spine(as))
--   .lambda b   ->  pair("lambda", enc(b))
--   etc.
\end{lstlisting}

\code{patternToLPTerm} is mutually recursive with \code{patternListToLPSpine}
(right-spine encoding of lists).  The \emph{ground} version
\code{patternToLPGroundTerm} treats \code{.fvar x} as \code{GroundTerm.const "fv:x"}
(making the variable name part of the constant).

The \emph{fragment} predicate \code{morkTranslatable} (inherited from the MORK
bridge) excludes patterns with \code{.subst} and rest-variable
\code{.collection} (i.e., \code{.collection ct elems (some rv)}):

\begin{lstlisting}[language=lean]
-- from MORK bridge:
def morkTranslatable : Pattern -> Bool
  | .fvar _              => true
  | .bvar _              => true
  | .apply _ args        => morkTranslatableList args
  | .lambda body         => morkTranslatable body
  | .multiLambda _ body  => morkTranslatable body
  | .collection _ _ none => morkTranslatableList elems
  | .collection _ _ _    => false   -- rest variable: excluded
  | .subst _ _           => false   -- explicit subst: excluded
\end{lstlisting}

\subsection{Commutation Lemma}

The key algebraic property is that grounding and encoding commute:

\begin{lemma}[\code{groundTerm\_commutes\_applyBindings}]
\label{lem:commute}
For \code{morkTranslatable} patterns $p$ and bindings $bs$:
\[
  \begin{aligned}
  &\code{bindingsToGrounding}(bs).\code{groundTerm}( \\
  &\qquad \code{patternToLPTerm}(p)) \\
  &= \code{patternToLPGroundTerm}( \\
  &\qquad \code{applyBindings}(bs, p))
  \end{aligned}
\]
\end{lemma}

\begin{proof}
By mutual well-founded recursion on \code{sizeOf p} and \code{sizeOf ps}
(for the list version).  The interesting cases:
\begin{itemize}
\item \code{.fvar x}: if $bs$ contains $(x, v)$, both sides evaluate to
  \code{patternToLPGroundTerm v}; otherwise both give \code{GroundTerm.const "fv:x"}.
\item \code{.bvar n}: both sides give \code{GroundTerm.const "bv:n"} (const on both
  sides by construction).
\item \code{.apply c args}: uses \code{groundTerm\_pairTerm} ($@[simp]$) and the
  induction hypothesis for the spine.
\item \code{.subst}, \code{.collection \_ \_ (some \_)}: excluded by
  \code{morkTranslatable}.
\end{itemize}
The helper \code{@[simp]} lemma \code{groundTerm\_pairTerm} establishes
that grounding distributes over \code{pairTerm} via \code{Matrix.cons\_val\_zero}
and \code{Matrix.cons\_val\_one}.
\end{proof}

\subsection{Top-Rule Completeness}

Rewrite rules are compiled to LP clauses with empty bodies:

\begin{lstlisting}[language=lean]
def rewriteRuleToLPClause (r : RewriteRule) : Clause mettailLPSig where
  head.symbol := ()
  head.args   := ![patternToLPTerm r.left, patternToLPTerm r.right]
  body        := []
\end{lstlisting}

\begin{theorem}[\code{lp\_complete\_topRule}]
\label{thm:topRule}
Let $r \in \code{lang.rewrites}$ be a rule with $r.\code{premises} = []$
and $\code{morkTranslatable}(r.left) = \code{true}$,
$\code{morkTranslatable}(r.right) = \code{true}$.
If $\code{applyBindings}(bs, r.left) = p$ and
$\code{applyBindings}(bs, r.right) = q$,
then
$\code{encodeReduces}(p, q) \in \mathrm{LHM}(\code{languageDefToLPKB}(\code{lang}))$.
\end{theorem}

\begin{proof}
(1) The LP clause $c = \code{rewriteRuleToLPClause}(r)$ is in \code{KB.prog}
(by \code{List.mem\_append} and \code{List.mem\_map}).
(2) The body of $c$ is empty, so the body condition is vacuously satisfied.
(3) The commutation lemma (Lemma~\ref{lem:commute}) gives:
\[
  g.\code{groundTerm}(\code{patternToLPTerm}(r.left)) = \code{patternToLPGroundTerm}(p)
\]
where $g = \code{bindingsToGrounding}(bs)$, and similarly for $r.right$ and $q$.
(4) A manual finitary argument over $\code{Fin 2}$ (using
\code{Matrix.cons\_val\_zero} and \code{Matrix.cons\_val\_one}) gives
$g.\code{groundAtom}(c.head) = \code{encodeReduces}(p, q)$.
(5) Theorem~\ref{thm:clause} (\code{leastHerbrandModel\_clause}) with steps (1)--(4)
gives the conclusion; using Theorem~\ref{thm:lfp} to convert from $T_P(LHM)$ to $LHM$.
\end{proof}

\subsection{Congruence Completeness}

For collection rewriting (the \code{congElem} case of \code{DeclReduces}),
two additional clause schemas are added to the LP:

\begin{lstlisting}[language=lean]
-- Head position:
-- reduces(coll(pair(x,xs)), coll(pair(y,xs))) :- reduces(x,y).
def collectionCongHeadClause (ct : CollType) : Clause mettailLPSig := ...

-- Tail position (recursive):
-- reduces(coll(pair(x,xs)), coll(pair(x,ys))) :- reduces(coll(xs),coll(ys)).
def collectionCongTailClause (ct : CollType) : Clause mettailLPSig := ...
\end{lstlisting}

\begin{lemma}[\code{lp\_complete\_congElem\_none}]
\label{lem:cong}
Assume $\code{encodeReduces}(\text{elems}[i], q') \in \mathrm{LHM}$,
$i < |\text{elems}|$, and \code{ct} is an enabled congruence collection kind.
Then replacing the $i$-th element by $q'$ preserves encoded reducibility:
\[
  \code{encodeReduces}(C, C[i := q']) \in \mathrm{LHM},
\]
where $C = \code{.collection ct elems none}$.
\end{lemma}

\begin{proof}
By list induction.  For the head position ($i = 0$): instantiate
\code{collectionCongHeadClause ct} with the grounding $g$ mapping
$\text{"x"} \mapsto \text{enc}(e)$, $\text{"y"} \mapsto \text{enc}(q')$,
$\text{"xs"} \mapsto \text{enc}(\text{rest})$.  The body atom reduces to
\code{encodeReduces}$(e, q')$, which is in the model by hypothesis.
Theorem~\ref{thm:clause} gives the head.  For the tail position ($i = n+1$):
use the induction hypothesis on \code{rest}, then instantiate
\code{collectionCongTailClause ct} to lift.
\end{proof}

\subsection{Full Completeness}

\begin{theorem}[\code{lp\_complete}]
\label{thm:complete}
If $\code{DeclReduces}(\code{lang}, p, q)$,
all rules in \code{lang.rewrites} are \code{lpTranslatable},
and $\code{morkTranslatable}(p) = \code{true}$,
then $\code{encodeReduces}(p, q) \in \mathrm{LHM}(\code{languageDefToLPKB}(\code{lang}))$.
\end{theorem}

\begin{proof}
By cases on the \code{DeclReduces} derivation:
\begin{itemize}
\item \textbf{\code{topRule}}: The rule $r$ is \code{lpTranslatable}, so both sides are
  mork-translatable and the LHS is \code{isMatchCorrect}.  The hypothesis
  $bs \in \code{matchPattern}(r.left, p)$ combined with
  Theorem~\ref{thm:matchcorrect} (\code{matchPattern\_correct}) gives
  $\code{applyBindings}(bs, r.left) = p$.  Then Theorem~\ref{thm:topRule}
  (\code{lp\_complete\_topRule}) gives the conclusion.
\item \textbf{\code{congElem}}: The inner step is handled by \code{lp\_complete\_topRule};
  then Lemma~\ref{lem:cong} (\code{lp\_complete\_congElem\_none}) lifts to the
  collection context.
\end{itemize}
\end{proof}

The LP-translatability predicate is:
\begin{lstlisting}[language=lean]
def lpTranslatable (r : RewriteRule) : Bool :=
  morkTranslatable r.left && morkTranslatable r.right &&
  Pattern.isMatchCorrect r.left
\end{lstlisting}

%% ====================================================================
\section{\PeTTa{} Evaluation and LP Soundness}
\label{sec:petta}
%% ====================================================================

\subsection{\PeTTa{} Atomspace}

\code{OSLF/PeTTa/SpaceSemantics.lean} (160 lines) defines the atomspace structure:

\begin{lstlisting}[language=lean]
structure PeTTaSpace where
  facts : List Pattern    -- EDB ground atoms
  rules : List RewriteRule

namespace PeTTaSpace
  def empty   : PeTTaSpace := { facts := [], rules := [] }
  def addAtom (s : PeTTaSpace) (p : Pattern) : PeTTaSpace :=
    { s with facts := p :: s.facts }
  def removeAtom (s : PeTTaSpace) (p : Pattern) : PeTTaSpace :=
    { s with facts := s.facts.filter (fun x => not (x == p)) }
  def addRule (s : PeTTaSpace) (r : RewriteRule) : PeTTaSpace :=
    { s with rules := r :: s.rules }
\end{lstlisting}

The space query operation models \code{(match \&self pat tmpl)} in \MeTTa{}:

\begin{lstlisting}[language=lean]
def spaceMatch (s : PeTTaSpace) (pat tmpl : Pattern) : Answers :=
  s.facts.flatMap fun fact =>
    (matchPattern pat fact).map fun bs => applyBindings bs tmpl
\end{lstlisting}

Soundness and completeness of \code{spaceMatch} are proven:
every answer arises from a fact-match pair, and every fact-match pair
produces an answer.

\subsection{\PeTTa{} Evaluation Relation}

\code{OSLF/PeTTa/Eval.lean} (187 lines) defines the inductive evaluation judgment:

\begin{lstlisting}[language=lean]
inductive PeTTaEval (s : PeTTaSpace) : Pattern -> Answers -> Prop where

  | var (x : String) :
      PeTTaEval s (.fvar x) [.fvar x]

  | bvar (n : Nat) :
      PeTTaEval s (.bvar n) [.bvar n]

  | ground (c : String) :
      PeTTaEval s (.apply c []) [.apply c []]

  | ruleApp (r : RewriteRule) (bs : Bindings) (p q : Pattern)
      (hr    : r in s.rules)
      (hprem : r.premises = [])
      (hm    : bs in matchPattern r.left p)
      (hq    : applyBindings bs r.right = q) :
      PeTTaEval s p [q]

  | spaceQuery (pat tmpl : Pattern) (results : Answers)
      (hres : results = s.spaceMatch pat tmpl) :
      PeTTaEval s (.apply "match" [.apply "&self" [], pat, tmpl]) results

  | superpose (alts : List Pattern) :
      PeTTaEval s (.apply "superpose" [.collection .vec alts none]) alts

  | collapse (p : Pattern) (answers : Answers)
      (h : PeTTaEval s p answers) :
      PeTTaEval s (.apply "collapse" [p]) [.collection .vec answers none]
\end{lstlisting}

The alignment with the HE MeTTa specification and the PeTTa transpiler is:

\begin{center}
\begin{tabular}{lll}
\hline
\textbf{MeTTa spec} & \textbf{PeTTaEval} & \textbf{PeTTa (Prolog)} \\
\hline
\code{metta(Variable, ...)} & \code{var} & clause head for fvar \\
\code{metta(Atom, ...)} & \code{ground} & fact clause \\
\code{metta\_call(rhs, ...)} & \code{ruleApp} & top-level rule clause \\
\code{metta(match \&self ...)} & \code{spaceQuery} & \code{match\_term/3} \\
\code{superpose-bind} & \code{superpose} & Prolog \code{;} disjunction \\
\code{collapse-bind} & \code{collapse} & \code{findall/3} \\
\hline
\end{tabular}
\end{center}

Key properties proven in \code{Eval.lean}:
\begin{itemize}
\item \code{petta\_eval\_var\_case}: \code{PeTTaEval s (.fvar x) [.fvar x]} holds for all $x$.
\item \code{petta\_eval\_ruleApp\_singleton}: rule application always produces exactly one answer.
\item \code{petta\_eval\_superpose\_nil}: \code{PeTTaEval s (superpose []) []}.
\item \code{petta\_eval\_collapse\_singleton}: \code{collapse} always produces a singleton.
\item \code{spaceMatch\_mono\_addAtom}: adding a fact to the space only extends \code{spaceMatch} answers.
\end{itemize}

\PeTTa{} is inherently non-deterministic: for a given expression $p$, multiple
constructors of \code{PeTTaEval} may fire simultaneously (e.g., \code{.fvar x}
can match both the \code{var} constructor and a \code{ruleApp} with an
fvar-matching LHS).

\subsection{Compiling \PeTTa{} Spaces to LP}

\code{OSLF/PeTTa/LPSoundness.lean} (255 lines) defines the compilation:

\begin{lstlisting}[language=lean]
def pettaSpaceToLangDef (s : PeTTaSpace) : LanguageDef where
  rewrites              := s.rules
  congruenceCollections := []   -- flat rules only
  ...

def pettaSpaceToLPKB (s : PeTTaSpace) : KnowledgeBase mettailLPSig :=
  languageDefToLPKB (pettaSpaceToLangDef s)
\end{lstlisting}

Only \code{s.rules} is compiled; \code{s.facts} are handled by \code{spaceMatch}.

\subsection{LP Soundness}

\begin{definition}[\code{pettaRuleSafe}]
A rule $r$ is \emph{PeTTa LP-safe} if:
$\code{morkTranslatable}(r.left) = \code{true}$,
$\code{morkTranslatable}(r.right) = \code{true}$, and
$\code{r.left.isMatchCorrect} = \code{true}$.
This is exactly \code{lpTranslatable r}.
\end{definition}

\begin{theorem}[\code{petta\_ruleApp\_lp\_sound}]
\label{thm:sound}
Let $r$ be LP-safe with $r \in s.\code{rules}$, $r.\code{premises} = []$,
$bs \in \code{matchPattern}(r.left, p)$, $\code{applyBindings}(bs, r.right) = q$.
Then $\code{encodeReduces}(p, q) \in \mathrm{LHM}(\code{pettaSpaceToLPKB}(s))$.
\end{theorem}

\begin{proof}
(1) \code{pettaRuleSafe\_iff} unpacks the safety into the three components.
(2) Theorem~\ref{thm:matchcorrect} (\code{matchPattern\_correct}) applied to
$bs \in \code{matchPattern}(r.left, p)$ with \code{isMatchCorrect} gives
$\code{applyBindings}(bs, r.left) = p$.
(3) $r \in s.\code{rules}$ lifts to $r \in (\code{pettaSpaceToLangDef}(s)).\code{rewrites}$ by \code{rfl}.
(4) Theorem~\ref{thm:topRule} (\code{lp\_complete\_topRule}) gives the conclusion.
\end{proof}

\begin{corollary}[\code{petta\_safe\_space\_ruleApp\_lp\_sound}]
\label{cor:safe}
For an LP-safe space $s$ (all rules satisfy \code{pettaRuleSafe}),
every \code{ruleApp} derivation $r \in s.\code{rules}$, $bs$, $p$, $q$
is witnessed in $\mathrm{LHM}(\code{pettaSpaceToLPKB}(s))$.
\end{corollary}

\begin{proof}
Immediate from Theorem~\ref{thm:sound} using $hs(r, hr)$ for the safety condition.
\end{proof}

\begin{theorem}[\code{pettaSpaceToLPKB\_addRule\_mono}]
Adding a rule to a space preserves LP model membership: if
$a \in \mathrm{LHM}(\code{pettaSpaceToLPKB}(s))$, then
$a \in \mathrm{LHM}(\code{pettaSpaceToLPKB}(s.\code{addRule}(r)))$.
\end{theorem}

\begin{proof}
Key sub-lemma: for any interpretation $I$,
$T_P(\code{KB}(s), I) \subseteq T_P(\code{KB}(s.\code{addRule}(r)), I)$
(proven by noting that $s.\code{rules} \subseteq r :: s.\code{rules}$ and
using \code{List.mem\_cons\_of\_mem}).
Then \code{leastHerbrandModel\_least} (Theorem~\ref{thm:least}) with
$I = \mathrm{LHM}(s.\code{addRule}(r))$ being a pre-fixpoint of $T_P(KB(s))$
gives $\mathrm{LHM}(s) \subseteq \mathrm{LHM}(s.\code{addRule}(r))$.
\end{proof}

A sanity check completes the file: the empty space compiles to the empty KB,
$\code{pettaSpaceToLPKB}(\code{empty}) = \{prog := [], db := \emptyset\}$.

\subsection{Unification}

The files \code{Unification.lean} (259 lines) and \code{UnificationMGU.lean}
(229 lines) formalise unification for the LP kernel.  The most general unifier (MGU)
is computed and proved correct: \code{unify\_mgu} states that any unifier factors
through the MGU.  These results are used in the SLD resolution soundness proof
but are stated over the full LP signature (not restricted to \code{mettailLPSig}).

%% ====================================================================
\section{Certified Propositional ATP/Chainer Layer}
\label{sec:chainer}
%% ====================================================================

\code{Logic/LP/PropositionalConnectionChainer.lean} (742 lines) provides a
connection-tableau proof object language aligned with the \PeTTa{} leanCoP-style
canaries.  The module contains:
\begin{itemize}
\item \code{TraceDerivable}/\code{ConnDerivable} (abstract derivability),
\item executable replay checking (\code{replay}, \code{checkProof}),
\item executable DFS search (\code{connProveDFS}),
\item refinement and witness-level completeness theorems.
\end{itemize}

The main executable refinement theorems are:

\begin{theorem}[\code{connProveDFS\_sound}]
\label{thm:connsound}
If \code{connProveDFS} returns \code{some tr}, then the goal is abstractly
connection-derivable.
\end{theorem}

\begin{theorem}[\code{connProveDFS\_witness\_complete}]
\label{thm:conncomplete}
If a goal is connection-derivable, then there exist finite fuel and a returned
trace from \code{connProveDFS}.
\end{theorem}

The fixture map section \code{PeTTaFixtureMap} provides theorem-level parity with
\PeTTa{} \code{ic\_*} canaries.

\begin{table}[H]
\centering
\begin{tabularx}{\linewidth}{YYY}
\hline
\textbf{PeTTa canary intent} & \textbf{Lean theorem} & \textbf{Expected outcome} \\
\hline
UNSAT proof object accepted &
\codepath{ic_unsat_checkProof_true} &
\code{checkProof = true} \\
SAT singleton rejects DFS proof &
\codepath{ic_sat_connProve_none} &
\code{connProveDFS = none} (all fuel) \\
UNSAT DFS finds witness &
\codepath{ic_unsat_connProve_has_proof} &
$\exists tr,\ \code{some}\ tr$ at fuel 2 \\
Branching proof (root $\neg p$) accepted &
\codepath{ic_branching_negp_checkProof_true} &
\code{checkProof = true} \\
Branching proof (root $\neg q$) accepted &
\codepath{ic_branching_negq_checkProof_true} &
\code{checkProof = true} \\
Two distinct branching roots both valid &
\codepath{ic_branching_two_distinct_root_proofs} &
root inequality + both checks true \\
Branching DFS witness ($\neg p$ root) &
\codepath{ic_branching_negp_connProve_has_proof} &
$\exists tr,\ \code{some}\ tr$ at fuel 3 \\
Branching DFS witness ($\neg q$ root) &
\codepath{ic_branching_negq_connProve_has_proof} &
$\exists tr,\ \code{some}\ tr$ at fuel 3 \\
\hline
\end{tabularx}
\caption{Formal theorem map for \PeTTa{} \code{ic\_*} propositional canaries.}
\label{tab:icmap}
\end{table}

FO status is tracked in \code{Logic/LP/FirstOrderConnectionTrace.lean}.  The
proven FO endpoints now include replay-checker soundness
(\code{replayFO\_sound}), executable DFS soundness
(\code{connProveFODFS\_sound}), and witness-level DFS completeness
(\code{connProveFODFS\_witness\_complete}).

This section therefore establishes theorem-prover-style executable refinement in
Lean for the propositional fragment, with explicit FO soundness scaffolding already
formalized and compiled.

%% ====================================================================
\section{File Statistics}
\label{sec:stats}
%% ====================================================================

\begin{table}[H]
\centering
\scriptsize
\begin{tabularx}{\linewidth}{Y r Y c}
\hline
\textbf{File} & \textbf{Lines} & \textbf{Key theorems} & \textbf{Sorries} \\
\hline
\multicolumn{4}{l}{\textit{LP Kernel (\code{Logic/LP/})}} \\
\code{Core.lean}              & 243 & LP signature, terms, clauses & 0 \\
\code{Substitution.lean}      & 206 & grounding algebra & 0 \\
\code{Semantics.lean}         & 204 & $T_P$, LHM, fixpoints & 0 \\
\code{Matching.lean}          & 333 & LP term matching & 0 \\
\code{SLD.lean}               & 169 & SLD soundness & 0 \\
\code{SLDAll.lean}            & 205 & all-solutions SLD & 0 \\
\code{SLDCompute.lean}        & 485 & semantic-lift completeness & 0 \\
\code{SLDCompletenessKit.lean} & 261 & canonical witness kits & 0 \\
\code{SLDCompletenessCanaries.lean} & 178 & completeness canaries & 0 \\
\code{Unification.lean}       & 259 & unifyFuel soundness & 0 \\
\code{UnificationMGU.lean}    & 229 & MGU factoring & 0 \\
\code{UnificationComplete.lean} & 513 & \codepath{unifyFuel_exists_of_unifies} & 0 \\
\code{UnificationCompletenessCanaries.lean} & 122 & occurs/non-occurs canaries & 0 \\
\code{PropositionalConnectionChainer.lean} & 1{,}245 & DFS soundness/completeness & 0 \\
\code{PropositionalChainer.lean} & 265 & propositional chainer & 0 \\
\code{FirstOrderConnectionTrace.lean} & 1{,}440 & FO connection trace & 0 \\
\code{MeTTaILBridge.lean}     & 847 & commutation, LP completeness & 0 \\
\code{FunctionFree.lean}      & 109 & function-free fragment & 0 \\
\code{FunctionFreeEvaluation.lean} & 110 & FF evaluation & 0 \\
\code{FunctionFreePeTTa.lean} & 481 & FF PeTTa integration & 0 \\
\code{Provenance.lean}        & 125 & proof provenance & 0 \\
\code{RangeRestriction.lean}  & 197 & range restriction & 0 \\
\code{MMMeasure.lean}         & 344 & MM measure theory & 0 \\
\code{OSLFBridge.lean}        & 172 & LP--OSLF bridge & 0 \\
\code{PathMapBridge.lean}     & 97  & LP--PathMap bridge & 0 \\
\code{WorldModelBridge.lean}  & 92  & LP--WM bridge & 0 \\
\code{CertifyingDatalogBridge.lean} & 82 & Datalog bridge & 0 \\
\hline
\multicolumn{4}{l}{\textit{MeTTaIL Matching (\code{OSLF/MeTTaIL/})}} \\
\code{Match.lean}             & 201 & algorithmic matcher & 0 \\
\code{MatchSpec.lean}         & 821 & \codepath{matchPattern_correct} & 0 \\
\hline
\multicolumn{4}{l}{\textit{PeTTa Specification (\code{OSLF/PeTTa/})}} \\
\code{Answers.lean}           & 121 & answer type & 0 \\
\code{SpaceSemantics.lean}    & 160 & \codepath{spaceMatch} & 0 \\
\code{Eval.lean}              & 187 & \codepath{PeTTaEval} & 0 \\
\code{LPSoundness.lean}       & 255 & ruleApp LP soundness & 0 \\
\code{Effects.lean}           & 301 & \code{PeTTaCmd} + state monotonicity & 0 \\
\code{TypeSystem.lean}        & 374 & \code{MeTTaType} arrow typing & 0 \\
\code{PrologBridge.lean}      & 125 & Prolog goal compilation & 0 \\
\code{TranslateExpr.lean}     & 463 & \codepath{compileExpr} correctness & 0 \\
\code{MeTTaEval.lean}         & 398 & 4-arg \code{MeTTaEval} judgment & 0 \\
\code{StdLib.lean}            & 309 & \code{if-equal}, \code{car-atom}, \code{let} & 0 \\
\code{GroundedOracle.lean}    & 297 & grounded function oracles & 0 \\
\code{TypedEval.lean}         & 296 & type-driven evaluation & 0 \\
\code{MinimalInstructions.lean} & 388 & minimal instruction set & 0 \\
\code{OSLFInstance.lean}      & 144 & \codepath{pettaOSLF}, Galois, Rust export & 0 \\
\code{GSLTVertex.lean}        & 103 & GSLT forward fiber & 0 \\
\hline
\textbf{Total (44 files)} & \textbf{13{,}956} & & \textbf{0} \\
\hline
\end{tabularx}
\normalsize
\caption{File statistics for the certified \PeTTa{} formalization.  All files use only
  \code{[propext, Classical.choice, Quot.sound]}.}
\label{tab:stats}
\end{table}

%% ====================================================================
\section{Stateful Evaluation: Effects and Commands}
\label{sec:effects}
%% ====================================================================

\code{OSLF/PeTTa/Effects.lean} (277 lines) extends the pure \code{PeTTaEval}
relation with a stateful layer that models \PeTTa{}'s mutating atomspace
operations.

\subsection{Motivation}

The pure \code{PeTTaEval} relation (Section~\ref{sec:petta}) captures rule
application, space queries, \code{superpose}, and \code{collapse}, but it treats
the atomspace as a static read-only resource.  The real \PeTTa{} interpreter
supports imperative mutations: \code{add-atom} and \code{remove-atom} modify
\code{\&self} in place; \code{get-atoms} retrieves all current facts; and
\code{progn}/\code{prog1} sequence effects in a specific return-value discipline.
Formalizing this requires a stateful layer.

\subsection{Evaluation State}

\begin{lstlisting}[language=lean]
structure EvalState where
  space : PeTTaSpace
\end{lstlisting}

\code{EvalState} wraps a \code{PeTTaSpace} (the single \code{\&self} atomspace).
The namespace provides five operations:

\begin{lstlisting}[language=lean]
def EvalState.empty     : EvalState
def EvalState.addAtom   (s : EvalState) (p : Pattern) : EvalState
def EvalState.removeAtom (s : EvalState) (p : Pattern) : EvalState
def EvalState.addRule   (s : EvalState) (r : RewriteRule) : EvalState
def EvalState.withSpace (s : EvalState) (sp : PeTTaSpace) : EvalState
\end{lstlisting}

All operations are purely functional: they return a new \code{EvalState}.
The unit return value for \code{add-atom} and \code{remove-atom} is:

\begin{lstlisting}[language=lean]
def unitAtom : Pattern := .apply "()" []
\end{lstlisting}

\subsection{The \texttt{PeTTaCmd} Judgment}

The stateful judgment \code{PeTTaCmd s\textsubscript{0} expr s\textsubscript{1} answers}
means: starting from state \code{s\textsubscript{0}}, evaluating \code{expr}
transitions to state \code{s\textsubscript{1}} and produces nondeterministic
answer set \code{answers}.

\begin{lstlisting}[language=lean]
inductive PeTTaCmd : EvalState -> Pattern -> EvalState -> Answers -> Prop where
  | addAtomCmd (s : EvalState) (p : Pattern) :
      PeTTaCmd s
        (.apply "add-atom" [.apply "&self" [], p])
        (s.addAtom p)
        [unitAtom]
  | removeAtomCmd (s : EvalState) (p : Pattern) :
      PeTTaCmd s
        (.apply "remove-atom" [.apply "&self" [], p])
        (s.removeAtom p)
        [unitAtom]
  | getAtomsCmd (s : EvalState) :
      PeTTaCmd s
        (.apply "get-atoms" [.apply "&self" []])
        s
        s.space.facts
  | pureEval (s : EvalState) (p : Pattern) (answers : Answers)
      (h : PeTTaEval s.space p answers) :
      PeTTaCmd s p s answers
  | prognCmd (s0 s1 s2 : EvalState) (e1 e2 : Pattern) (ans1 ans2 : Answers)
      (h1 : PeTTaCmd s0 e1 s1 ans1) (h2 : PeTTaCmd s1 e2 s2 ans2) :
      PeTTaCmd s0 (.apply "progn" [e1, e2]) s2 ans2
  | prog1Cmd (s0 s1 s2 : EvalState) (e1 e2 : Pattern) (ans1 ans2 : Answers)
      (h1 : PeTTaCmd s0 e1 s1 ans1) (h2 : PeTTaCmd s1 e2 s2 ans2) :
      PeTTaCmd s0 (.apply "prog1" [e1, e2]) s2 ans1
\end{lstlisting}

The six constructors cover all effectful and pure evaluation forms.
The \code{pureEval} constructor embeds any \code{PeTTaEval} derivation
state-preservingly into the stateful layer.  The \code{prognCmd} constructor
sequences two commands, threading the intermediate state \code{s\textsubscript{1}}:
the second command is evaluated in the output state of the first, and the
final answers are those of the second.  \code{prog1Cmd} is symmetric but
returns the answers of the first command.

\subsection{Correspondence with PeTTa Prolog}

\begin{table}[H]
\centering
\begin{tabular}{ll}
\hline
\textbf{PeTTa Prolog predicate} & \textbf{\code{PeTTaCmd} constructor} \\
\hline
\code{add\_atom\_to\_space/2}                       & \code{addAtomCmd}   \\
\code{remove\_atom\_from\_space/2}                  & \code{removeAtomCmd} \\
\code{findall(A, get\_atom(self,A), As)}             & \code{getAtomsCmd}  \\
pure \code{metta\_call/2}                           & \code{pureEval}     \\
\code{progn/2} (from \code{lib\_metta4})            & \code{prognCmd}     \\
\code{prog1/2} (from \code{lib\_metta4})            & \code{prog1Cmd}     \\
\hline
\end{tabular}
\caption{Correspondence between PeTTa Prolog predicates and \code{PeTTaCmd} constructors.}
\label{tab:cmd-prolog}
\end{table}

\subsection{Key Theorems}

\begin{lemma}[\code{pureEval\_lifts}]
Any \code{PeTTaEval} derivation lifts to a state-preserving \code{PeTTaCmd}:
if \code{PeTTaEval s.space p ans}, then \code{PeTTaCmd s p s ans}.
\end{lemma}

\begin{lemma}[\code{addAtomCmd\_mem\_facts}]
The added atom is a member of the output state's fact list:
$p \in (\code{s.addAtom}\ p)\code{.space.facts}$.
\end{lemma}

\begin{lemma}[\code{addAtomCmd\_preserves\_facts}]
Existing facts survive an \code{addAtomCmd}: if $q \in \code{s.space.facts}$
then $q \in (\code{s.addAtom}\ p)\code{.space.facts}$.
\end{lemma}

\begin{lemma}[\code{removeAtomCmd\_subset\_facts}]
\code{removeAtomCmd} only removes the targeted atom: any fact surviving in
the output state was already present in the input state.
\end{lemma}

\begin{lemma}[\code{pettaCmd\_shape}]
Exhaustive case analysis: given any \code{PeTTaCmd s p s\textsubscript{1} ans},
exactly one of the following characterizes the expression form and state transition:
\begin{enumerate}
  \item $p = \code{(add-atom \&self}\ q\code{)}$ and $s_1 = \code{s.addAtom}\ q$,
  \item $p = \code{(remove-atom \&self}\ q\code{)}$ and $s_1 = \code{s.removeAtom}\ q$,
  \item $p = \code{(get-atoms \&self)}$ and $s_1 = s$ and $ans = \code{s.space.facts}$,
  \item $s_1 = s$ and $\code{PeTTaEval}\ \code{s.space}\ p\ ans$,
  \item $\exists\, e_1, e_2,\ p = \code{(progn}\ e_1\ e_2\code{)}$, or
  \item $\exists\, e_1, e_2,\ p = \code{(prog1}\ e_1\ e_2\code{)}$.
\end{enumerate}
\end{lemma}

\begin{theorem}[\code{example\_addThenGet}]
Concrete derivation: evaluating \code{(progn (add-atom \&self (foo)) (get-atoms \&self))}
from the empty state yields \code{[.apply "foo" []]}:
\begin{lstlisting}[language=lean]
theorem example_addThenGet :
    PeTTaCmd EvalState.empty
      (.apply "progn"
        [ .apply "add-atom" [.apply "&self" [], .apply "foo" []]
        , .apply "get-atoms" [.apply "&self" []] ])
      { space := { facts := [.apply "foo" []], rules := [] } }
      [.apply "foo" []] :=
  PeTTaCmd.prognCmd _ _ _ _ _ _ _
    (PeTTaCmd.addAtomCmd EvalState.empty (.apply "foo" []))
    (PeTTaCmd.getAtomsCmd _)
\end{lstlisting}
\end{theorem}

%% ====================================================================
\section{MeTTa Type System Fragment}
\label{sec:typesystem}
%% ====================================================================

\code{OSLF/PeTTa/TypeSystem.lean} (374 lines) formalizes the static fragment of
MeTTa's type system as an inductive judgment \code{MeTTaType s p t},
meaning ``in atomspace \code{s}, pattern \code{p} has type \code{t}''.

\subsection{Special Type Atoms}

MeTTa uses a handful of distinguished type-level atoms:

\begin{lstlisting}[language=lean]
def undefinedType   : Pattern := .apply "%Undefined%" []  -- supertype
def emptyType       : Pattern := .apply "%Empty%" []      -- bottom
def atomType        : Pattern := .apply "Atom" []         -- nullary apps
def expressionType  : Pattern := .apply "Expression" []   -- all apps
def arrowType (A B : Pattern) : Pattern := .apply "->" [A, B]
def typeAnnotationPat (p t : Pattern) : Pattern := .apply ":" [p, t]
\end{lstlisting}

Type annotations are stored as facts in the atomspace: \code{(: p t) $\in$ s.facts}
directly asserts that \code{p} has type \code{t}.

\subsection{The Type Judgment}

\begin{lstlisting}[language=lean]
inductive MeTTaType (s : PeTTaSpace) : Pattern -> Pattern -> Prop where

  | typeAnnotation (p t : Pattern)
      (h : typeAnnotationPat p t in s.facts) :
      MeTTaType s p t

  | undefinedIsTop (p : Pattern) :
      MeTTaType s p undefinedType

  | arrowApp (c : String) (x argTy retTy : Pattern)
      (hf : MeTTaType s (.apply c []) (arrowType argTy retTy))
      (hx : MeTTaType s x argTy) :
      MeTTaType s (.apply c [x]) retTy

  | arrowAppUnknown (c : String) (x argTy retTy : Pattern)
      (hf : MeTTaType s (.apply c []) (arrowType argTy retTy)) :
      MeTTaType s (.apply c [x]) undefinedType

  | arrowMultiApp (c : String) (x argTy midTy retTy : Pattern)
      (hf : MeTTaType s (.apply c []) (arrowType argTy (arrowType midTy retTy)))
      (hx : MeTTaType s x argTy) :
      MeTTaType s (.apply c [x]) (arrowType midTy retTy)

  | symbolIsAtom (c : String) :
      MeTTaType s (.apply c []) atomType

  | appIsExpression (c : String) (args : List Pattern) (hne : args /= []) :
      MeTTaType s (.apply c args) expressionType

  | varIsUndefined (x : String) :
      MeTTaType s (.fvar x) undefinedType
\end{lstlisting}

The eight constructors cover the key cases from the MeTTa spec's
\code{check\_if\_function\_type\_is\_applicable} and \code{check\_argument\_type}
predicates.  Arrow constructors use \code{c : String} as the function head
(since \code{Pattern.apply : String $\to$ List Pattern $\to$ Pattern} requires a
string head); a bare symbol \code{f} is represented as \code{.apply c []}.

\subsection{Design Decisions}

\paragraph{Types as patterns.}  The type language is \emph{the same} language as
the term language; there is no separate sort.  This mirrors MeTTa's philosophy of
a single, uniform atom space.

\paragraph{Non-determinism.}  \code{MeTTaType} is a relation, not a function:
a pattern may have multiple types simultaneously (e.g., every pattern has
\code{\%Undefined\%} via \code{undefinedIsTop}, and separately may have a
specific annotated type).

\paragraph{Non-dependent fragment.}  We do not model type-level computation,
dependent types, or the \code{Expression} type's special structural role in
the full MeTTa spec.  These are left for future work.

\paragraph{Function head restriction.}  Arrow rules require the function to be
a bare symbol (\code{.apply c []}), not an arbitrary expression.  The
\code{arrowApp} constructor therefore carries \code{c : String} rather than a
general pattern.  This is the normal use case in MeTTa (function symbols are
declared with \code{(: f (-> A B))} where \code{f} is a bare symbol).

\subsection{Key Theorems}

\begin{theorem}[\code{all\_spaces\_well\_typed}]
Every \code{PeTTaSpace} is well-typed: for all \code{p}, \code{t},
if \code{(: p t) $\in$ s.facts} then \code{MeTTaType s p t}.
\end{theorem}

\begin{proof}
Directly by \code{MeTTaType.typeAnnotation}.
\end{proof}

\begin{theorem}[\code{typeOf\_mono\_addAtom}]
Type judgments are monotone: if \code{MeTTaType s p t}, then
\code{MeTTaType (s.addAtom f) p t} for any new fact \code{f}.
\end{theorem}

\begin{proof}
By induction on the \code{MeTTaType} derivation, with
\code{PeTTaSpace.mem\_facts\_addAtom} for the annotation case.
\end{proof}

\begin{theorem}[\code{every\_pattern\_has\_undefined\_type}]
No pattern is type-less: \code{MeTTaType s p \%Undefined\%} for all \code{s} and \code{p}.
\end{theorem}

\begin{proof}
Immediate by \code{MeTTaType.undefinedIsTop}.
\end{proof}

\begin{theorem}[\code{arrowApp\_implies\_f\_has\_arrow\_type}]
If \code{(c x)} has a non-trivial type $t$ (i.e., $t \neq \code{\%Undefined\%}$,
$t \neq \code{Expression}$) that does not arise from a direct annotation
\code{(: (c x) t) $\in$ s.facts}, then the bare symbol \code{c} has some arrow type
$\code{(-> A B)}$ in \code{s}.
\end{theorem}

\begin{proof}
By cases on the \code{MeTTaType} derivation for \code{(.apply c [x])}; the
\code{typeAnnotation} case is ruled out by hypothesis, \code{undefinedIsTop}
by $t \neq \code{\%Undefined\%}$, \code{appIsExpression} by $t \neq \code{Expression}$;
the remaining arrow cases all carry an \code{hf} witness.
\end{proof}

\subsection{Annotated Types Query}

\begin{lstlisting}[language=lean]
def PeTTaSpace.annotatedTypes (s : PeTTaSpace) (p : Pattern) : List Pattern :=
  s.facts.filterMap fun fact =>
    match fact with
    | .apply ":" [p', t] => if p' == p then some t else none
    | _                  => none
\end{lstlisting}

\begin{theorem}[\code{annotatedTypes\_sound}]
Every type returned by \code{annotatedTypes s p} is derivable:
$t \in \code{annotatedTypes}\ s\ p \Rightarrow \code{MeTTaType}\ s\ p\ t$.
\end{theorem}

\begin{proof}
Via a private helper lemma \code{annotatedTypes\_filterMap\_aux} that performs
case analysis on the \code{filterMap} body: only \code{.apply ":" [p', t']}
facts where \code{p' = p} contribute, giving \code{(: p' t) $\in$ s.facts},
from which \code{typeAnnotation} applies.  The \code{c $\neq$ ":"} branch
uses \code{split} + head-symbol extraction to derive a contradiction.
\end{proof}

\subsection{Alignment with the MeTTa Specification}

\begin{center}
\begin{tabularx}{\linewidth}{YY}
\hline
\textbf{MeTTa spec predicate} & \textbf{\code{MeTTaType} constructor} \\
\hline
\code{(: p t) $\in$ space}                                   & \code{typeAnnotation} \\
\code{\%Undefined\%} as supertype                             & \code{undefinedIsTop} \\
\codepath{check_if_function_type_is_applicable}               & \code{arrowApp}, \code{arrowMultiApp} \\
\code{check\_argument\_type} (unknown arg)                    & \code{arrowAppUnknown} \\
Bare symbols have type \code{Atom}                            & \code{symbolIsAtom} \\
Non-nullary applications have type \code{Expression}          & \code{appIsExpression} \\
Unbound variables have type \code{\%Undefined\%}              & \code{varIsUndefined} \\
\hline
\end{tabularx}
\end{center}

The \code{arrowAppUnknown} constructor handles the MeTTa spec's fallback: when
argument type checking fails (the argument type does not match or is unknown),
the result type is \code{\%Undefined\%} rather than a type error.

%% ====================================================================
\section{Prolog Goal Language and Expression Compilation}
\label{sec:prolog}
%% ====================================================================

\PeTTa{}'s \code{translator.pl} compiles \MeTTa{} expressions into Prolog
goals before evaluation.  We formalize this pipeline in three files.

\subsection{Prolog Goal Language}

\code{Logic/Prolog/Prolog.lean} (referenced via \code{PrologBridge.lean}, 125 lines)
defines a Prolog-style goal language:

\begin{lstlisting}[language=lean]
inductive PrologGoal where
  | reduceCall  : List Pattern -> PrologGoal
  | spaceMatch  : Pattern -> Pattern -> PrologGoal
  | findall     : String -> PrologGoal -> PrologGoal
  | assertGoal  : Pattern -> PrologGoal
  | conjGoal    : PrologGoal -> PrologGoal -> PrologGoal
  | disjGoal    : PrologGoal -> PrologGoal -> PrologGoal
  | trueGoal    : PrologGoal
\end{lstlisting}

The \code{reduceCall} constructor models PeTTa's \code{reduce/2} predicate;
\code{findall} models \code{findall/3}; \code{assertGoal} models Prolog
\code{assert/1}.  The constructors mirror the Prolog goal forms emitted by
\code{translate\_expr/3}.

\subsection{MeTTa Prolog Oracle}

\code{PrologBridge.lean} wires the abstract \code{EvalOracle} to
concrete \PeTTa{} semantics:

\begin{lstlisting}[language=lean]
def meTTaPrologOracle (s : PeTTaSpace) : EvalOracle where
  call args outs := match args with
    | [e] => PeTTaEval s e outs
    | _   => False
\end{lstlisting}

The semantic contract: \code{oracle.call [e] outs $\leftrightarrow$
PeTTaEval s e outs}.

\subsection{Expression Compilation (\code{compileExpr})}

\code{TranslateExpr.lean} (463 lines) defines the compilation pass and proves
its correctness:

\begin{lstlisting}[language=lean]
def compileExpr : Pattern -> PrologGoal
  | .fvar x               => .reduceCall [.fvar x]
  | .apply c []           => .reduceCall [.apply c []]
  | .apply "superpose" [.collection .vec alts none] =>
      disjList (alts.map compileExpr)
  | .apply "collapse" [body] =>
      .findall "Out" (compileExpr body)
  | .apply "match" [.apply "&self" [], pat, tmpl] =>
      .spaceMatch pat tmpl
  | e                     => .reduceCall [e]
\end{lstlisting}

\begin{theorem}[\code{compileExpr\_ruleApp\_iff}]
For rule applications, compilation is sound and complete: a \code{reduceCall}
Prolog evaluation of \code{compileExpr (.apply c args)} produces answer \code{q}
if and only if \code{PeTTaEval s (.apply c args) [q]} holds via \code{ruleApp}.
\end{theorem}

\begin{theorem}[\code{compileExpr\_match\_correct}]
For space queries, \code{compileExpr (match \&self pat tmpl)} evaluates under
the Prolog oracle to exactly the answers produced by \code{PeTTaEval.spaceQuery}.
\end{theorem}

\begin{theorem}[\code{compileExpr\_superpose\_sound}]
For superpose, the compiled disjunction is sound: each branch of the
Prolog disjunction corresponds to a member of the superpose alternatives list.
\end{theorem}

The compilation covers \code{fvar}, ground atoms, \code{superpose},
\code{collapse}, \code{match}, \code{if}, \code{let}, \code{case}, and a
conservative catch-all for unrecognised forms.

%% ====================================================================
\section{4-Argument \MeTTa{} Evaluation}
\label{sec:mettaeval}
%% ====================================================================

\code{OSLF/PeTTa/MeTTaEval.lean} (398 lines) refines \code{PeTTaEval} with three
orthogonal features: binding threading, error propagation, and type-driven pass-through.

\subsection{The \code{MeTTaEval} Judgment}

\begin{lstlisting}[language=lean]
inductive MeTTaEval (s : PeTTaSpace)
    : Pattern -> Pattern -> Bindings -> EvalResult -> Prop where
  | symbolPassThrough (c : String) (ty : Pattern) (bs : Bindings) :
      MeTTaEval s (.apply c []) ty bs [(.apply c [], bs)]
  | varPassThrough (x : String) (ty : Pattern) (bs : Bindings) :
      MeTTaEval s (.fvar x) ty bs [(.fvar x, bs)]
  | errorPassThrough (atom msg : Pattern) (ty : Pattern) (bs : Bindings) :
      MeTTaEval s (.apply "Error" [atom, msg]) ty bs
        [(.apply "Error" [atom, msg], bs)]
  | ruleApp (r : RewriteRule) (mbs : Bindings) (p q : Pattern)
      (ty : Pattern) (bs : Bindings)
      (hr : r in s.rules) (hprem : r.premises = [])
      (hm : mbs in matchPattern r.left p)
      (hq : applyBindings mbs r.right = q) :
      MeTTaEval s p ty bs [(q, mbs ++ bs)]
  | spaceQuery ...
  | superpose ...
  | collapse ...
\end{lstlisting}

The type parameter \code{ty} flows through but does not affect reduction in the
current fragment---it gates pass-through for special types
(\code{\%Undefined\%}, \code{Atom}, \code{Expression}).

\subsection{Embedding into \code{PeTTaEval}}

\begin{theorem}[\code{meTTaEval\_ruleApp\_to\_pettaEval}]
The \code{ruleApp} case of \code{MeTTaEval} projects faithfully to
\code{PeTTaEval.ruleApp}.
\end{theorem}

Similar projection theorems hold for \code{spaceQuery}, \code{superpose},
and \code{collapse}.

\subsection{Standard Library (\code{StdLib.lean}, 309 lines)}

\code{StdLib.lean} defines derived evaluation forms for \PeTTa{}'s standard
operations: \code{if-equal} (conditional on pattern equality), \code{car-atom}
(head extraction), \code{cdr-atom} (tail extraction), \code{let} and
\code{let*} (sequential binding), \code{case} (pattern matching), and
\code{sealed} (scope isolation).  Each is defined as a specialized
\code{MeTTaEval} step and proved sound against the base evaluation relation.

\subsection{Grounded Oracles (\code{GroundedOracle.lean}, 297 lines)}

\code{GroundedOracle.lean} formalizes the interface for grounded (external)
functions in \MeTTa{}: arithmetic, string operations, and type-checking predicates.
A \code{GroundedFn} is a deterministic function \code{List Pattern $\to$
Option (List Pattern)}; the \code{groundedApp} evaluation constructor fires
when the head symbol is a registered grounded function and the oracle succeeds.

\subsection{Typed Evaluation (\code{TypedEval.lean}, 296 lines)}

\code{TypedEval.lean} integrates \code{MeTTaType} (Section~\ref{sec:typesystem})
with \code{MeTTaEval}: the \code{TypedEvalStep} judgment combines a type
derivation with an evaluation step, ensuring that the input expression and
its type are consistent.  Key properties include type preservation under
evaluation and type soundness for the pass-through cases.

%% ====================================================================
\section{\OSLF{} Pipeline and GSLT Integration}
\label{sec:oslf}
%% ====================================================================

\code{OSLF/PeTTa/OSLFInstance.lean} (144 lines) and
\code{OSLF/PeTTa/GSLTVertex.lean} (103 lines) compose the \PeTTa{} specification
stack into the \OSLF{} framework.

\subsection{OSLF Type System}

The key definition composes \code{pettaSpaceToLangDef} (from
Section~\ref{sec:petta}) with the generic \OSLF{} type synthesis:

\begin{lstlisting}[language=lean]
def pettaOSLF (s : PeTTaSpace) :=
  langOSLF (pettaSpaceToLangDef s) "Expr"

theorem pettaGalois (s : PeTTaSpace) :
    GaloisConnection (langDiamond (pettaSpaceToLangDef s))
                     (langBox (pettaSpaceToLangDef s)) :=
  langGalois (pettaSpaceToLangDef s)
\end{lstlisting}

The Galois connection $\Diamond \dashv \Box$ is obtained automatically from
the generic \OSLF{} infrastructure.

\begin{theorem}[\code{pettaDiamond\_spec}]
$\Diamond\varphi(p) \leftrightarrow \exists q,\ \code{langReduces}\ p\ q \wedge \varphi(q)$.
\end{theorem}

\begin{theorem}[\code{pettaBox\_spec}]
$\Box\varphi(p) \leftrightarrow \forall q,\ \code{langReduces}\ q\ p \to \varphi(q)$.
\end{theorem}

\subsection{LP Soundness Bridge}

\begin{theorem}[\code{pettaOSLF\_lp\_sound}]
For LP-safe \PeTTa{} spaces, any \OSLF{} diamond witness via \code{ruleApp}
is backed by LP model membership.
\end{theorem}

This composes \code{matchPattern\_correct} (Section~\ref{sec:matchspec}) with
\code{lp\_complete\_topRule} (Section~\ref{sec:bridge}) through the \OSLF{} pipeline.

\subsection{Rust Export}

\begin{lstlisting}[language=lean]
def pettaRenderRust (s : PeTTaSpace) : String :=
  renderLanguage (pettaSpaceToLangDef s)

def pettaWriteRust (path : FilePath) (s : PeTTaSpace) : IO Unit :=
  writeLanguage path (pettaSpaceToLangDef s)
\end{lstlisting}

These produce \code{language! \{ ... \}} macro text consumable by the
\code{mettail-rust} runtime.

\subsection{GSLT Forward Fiber}

\code{GSLTVertex.lean} provides a unit-indexed GSLT forward fiber, following the
governance pattern from the companion \OSLF{} paper:

\begin{lstlisting}[language=lean]
def pettaIdMorphism (s : PeTTaSpace) :
    ForwardMorphism (pettaSpaceToLangDef s)
                    (pettaSpaceToLangDef s) where
  mapTerm := id
  forward_sim _ q hred := <q, .single hred, rfl>

def pettaForwardFiber (s : PeTTaSpace) : ForwardFiber Unit where
  lang  := fun _ => pettaSpaceToLangDef s
  morph := fun _ => pettaIdMorphism s
\end{lstlisting}

This places \PeTTa{} as a degenerate (constant) vertex in the GSLT hypercube,
enabling categorical composition with other language fibers (PathMap, Governance,
Temporal).

%% ====================================================================
\section{Related Work}
\label{sec:related}
%% ====================================================================

\paragraph{LP semantics.}
The Herbrand model semantics for LP was established by van Emden and
Kowalski~\cite{vanEmdenKowalski1976}, who proved the equivalence of operational
(SLD resolution) and declarative (least model) semantics for Horn clause programs.
Lloyd's monograph~\cite{Lloyd1987} remains the standard reference.  Our
formalization follows Chapter~2 of Lloyd directly, using \code{OrderHom.lfp}
from Mathlib to realise Tarski's theorem on the set lattice.

\paragraph{Lean 4 LP/logic formalization.}
The Mettapedia project contains a separate Datalog formalization
(formerly in \code{Mettapedia/Logic/Datalog/}, now consolidated into \code{Logic/LP/}),
which covered the Datalog fragment without function symbols.  The present
LP kernel generalises this to function symbols, enabling the pattern encoding.
We are not aware of other full-stack LP-to-pattern-matching soundness formalizations
in Lean~4.

\paragraph{\MeTTa{} and Hyperon.}
MeTTa and the Hyperon architecture are described in Goertzel et al.~\cite{goertzel2009pln}.
The PLN reasoning layer is formalized in other Mettapedia papers
(see \code{Mettapedia/Logic/PLNCore.lean} et al.).
The \OSLF{} framework and its modal type system are formalized separately
in the companion paper~\cite{leanOSLF2025}.

\paragraph{Certified pattern matching.}
Pattern matching correctness for a locally-nameless $\lambda$-calculus was studied
by Aydemir et al.~\cite{AydemarPOPL2008}.  Our \code{isMatchCorrect} fragment
corresponds to the ``linear'' fragment where each pattern variable occurs exactly
once, plus the exclusion of collection patterns whose bag-matching non-determinism
prevents uniqueness of the reconstructed term.

\paragraph{Prolog formalization.}
Formalizations of SLD resolution correctness are a classical topic; our treatment
follows Lloyd's presentation~\cite{Lloyd1987} mechanized in Lean~4's type theory.

\paragraph{Connection tableaux and leanCoP lineage.}
Connection-style calculi go back to Bibel's matrix/connection method
tradition~\cite{Bibel1987}.  leanCoP operationalized this line with a compact
Prolog implementation and practical connection-tableau search control
~\cite{Otten2010leanCoP}.  Our propositional chainer layer is not a full reproof
of leanCoP, but it follows the same proof-object shape (reduce/extend traces),
and, crucially, adds Lean-certified replay soundness and executable DFS
refinement theorems.

%% ====================================================================
\section{Conclusion and Future Work}
\label{sec:conclusion}
%% ====================================================================

We have presented a Lean~4 formalization of the \PeTTa{} language comprising
13{,}956 lines across 44 files with zero sorries and zero custom axioms.
The formalization establishes a certified end-to-end pipeline: from LP
foundations (Herbrand model, SLD resolution, unification completeness) through
\MeTTaIL{} pattern matching correctness and LP soundness for \PeTTa{}
rule application, to a Prolog-style goal language with expression compilation,
a full 4-argument evaluation judgment with grounded oracles and standard
library forms, and finally an \OSLF{} type system with automatic
$\Diamond \dashv \Box$ Galois connection and Rust code export.

\paragraph{Future work.}
\begin{itemize}
\item \textbf{Bridge converse direction.} The current bridge completeness result
  shows $\code{DeclReduces} \Rightarrow \mathrm{LHM}$; the converse
  ($\mathrm{LHM} \Rightarrow \code{DeclReduces}$) requires a constrained
  encoding discipline and a constructive lifting theorem from LP witnesses back
  into \MeTTaIL{} reductions.  Range restriction (\code{RangeRestriction.lean})
  is partially formalized.
\item \textbf{First-order connection search refinement.} The propositional
  connection-tableau layer is fully certified; extending executable refinement
  and witness-level completeness to the full first-order
  trace/substitution layer (\code{FirstOrderConnectionTrace.lean}) is
  ongoing work.
\item \textbf{Dependent types.} The current type system fragment covers
  annotations, arrow application, and pass-through types.  Formalizing
  type-level computation and dependent types would complete the alignment
  with the HE \MeTTa{} spec.
\item \textbf{Multi-space evaluation.} The current formalization uses a single
  \code{\&self} atomspace.  Extending to multi-space evaluation
  (\code{new-space}, \code{bind!}) would capture the full \MeTTa{} import
  and module system.
\end{itemize}

\appendix
\section{Theorem Index}

\begin{table}[H]
\centering
\scriptsize
\begin{tabularx}{\linewidth}{YYc}
\hline
\textbf{Layer} & \textbf{Canonical theorem} & \textbf{Status} \\
\hline
\multicolumn{3}{l}{\textit{LP Kernel}} \\
Unification (FO, semantic $\Rightarrow$ executable) &
\codepath{unifyFuel_exists_of_unifies} &
proved \\
Unification (function-free) &
\codepath{unifyFuel_exists_of_unifies_functionFree} &
proved \\
SLD query completeness from semantic lift &
\codepath{sldQuery_complete_of_semantic_lift} &
proved \\
SLD query complete+sound bundle &
\codepath{sldQuery_complete_sound_bundle_of_semantic_lift} &
proved \\
\hline
\multicolumn{3}{l}{\textit{ATP/Chainer}} \\
Propositional connection DFS soundness &
\codepath{connProveDFS_sound} &
proved \\
Propositional connection DFS witness completeness &
\codepath{connProveDFS_witness_complete} &
proved \\
FO connection DFS soundness &
\codepath{connProveFODFS_sound} &
proved \\
FO connection DFS witness completeness &
\codepath{connProveFODFS_witness_complete} &
proved \\
\hline
\multicolumn{3}{l}{\textit{MeTTaIL/PeTTa Bridge}} \\
Pattern matching correctness &
\codepath{matchPattern_correct} &
proved \\
MeTTaIL--LP completeness &
\codepath{lp_complete} &
proved \\
PeTTa ruleApp LP soundness &
\codepath{petta_ruleApp_lp_sound} &
proved \\
\hline
\multicolumn{3}{l}{\textit{Prolog Compilation}} \\
compileExpr ruleApp correctness &
\codepath{compileExpr_ruleApp_iff} &
proved \\
compileExpr match correctness &
\codepath{compileExpr_match_correct} &
proved \\
compileExpr superpose soundness &
\codepath{compileExpr_superpose_sound} &
proved \\
\hline
\multicolumn{3}{l}{\textit{OSLF Pipeline}} \\
PeTTa Galois connection &
\codepath{pettaGalois} &
proved \\
OSLF LP soundness bridge &
\codepath{pettaOSLF_lp_sound} &
proved \\
GSLT forward fiber &
\codepath{pettaForwardFiber} &
proved \\
\hline
\end{tabularx}
\normalsize
\caption{Compact theorem index for the certified \PeTTa{} formalization.}
\label{tab:completeness-index}
\end{table}

\bibliographystyle{plain}
\bibliography{references}

\end{document}
